\documentclass[11pt]{article}
\usepackage[margin=1in]{geometry}
\usepackage{amsmath,amssymb,amsthm,mathtools}
\usepackage{booktabs}
\usepackage{array}
\usepackage{enumitem}
\setlist[enumerate]{leftmargin=*,itemsep=0.4em,topsep=0.2em}
\setlist[itemize]{leftmargin=*,itemsep=0.3em,topsep=0.2em}

\newcommand{\R}{\mathbb{R}}
\newcommand{\E}{\mathbb{E}}
\newcommand{\Pp}{\mathbb{P}}
\newcommand{\1}{\mathbf{1}}
\newcommand{\xx}{\mathbf{x}}
\newcommand{\ww}{\mathbf{w}}
\newcommand{\bb}{\mathbf{b}}
\newcommand{\argmin}{\mathop{\mathrm{arg\,min}}}

\title{IE206 Advanced Question Bank: Detailed Solutions}
\author{}
\date{\today}

\begin{document}
\maketitle
\tableofcontents
\newpage

\section{Linear Regression}

\subsection*{Q1. Laplace Noise and MLE Objective}
Assume the model
\[
y_i = \xx_i^\top\beta + \varepsilon_i,\qquad \varepsilon_i\stackrel{i.i.d.}{\sim}\text{Laplace}(0,b).
\]
The Laplace density is
\[
f_{\varepsilon}(e)=\frac{1}{2b}\exp\!\left(-\frac{|e|}{b}\right).
\]
Hence the likelihood is
\[
L(\beta,b)=\prod_{i=1}^n \frac{1}{2b}\exp\!\left(-\frac{|y_i-\xx_i^\top\beta|}{b}\right),
\]
and the negative log-likelihood is
\[
-\log L(\beta,b)=n\log(2b)+\frac{1}{b}\sum_{i=1}^n |y_i-\xx_i^\top\beta|.
\]
For fixed $b$, MLE for $\beta$ is
\[
\hat\beta_{\text{MLE}}=\argmin_\beta\sum_{i=1}^n|y_i-\xx_i^\top\beta|,
\]
not least squares. So under Laplace noise, \emph{L1 regression} (LAD) is the correct MLE objective.

Robustness comparison:
\begin{itemize}
\item L2 loss penalizes residual $r$ as $r^2$; outliers get very large influence.
\item L1 loss penalizes residual as $|r|$; influence grows much slower.
\item Therefore Laplace-MLE/L1 is more robust than Gaussian-MLE/L2 to heavy tails/outliers.
\end{itemize}

\subsection*{Q2. Perfect Collinearity}
Let two features satisfy $\xx_{\cdot,2}=c\,\xx_{\cdot,1}$ exactly. Then columns of $X$ are linearly dependent, so $\operatorname{rank}(X)<p$ and
\[
X^\top X\ \text{is singular.}
\]

OLS objective:
\[
J(\beta)=\|y-X\beta\|_2^2.
\]
If $v\in\mathcal{N}(X)$ (null space), then $X(\beta+v)=X\beta$, so
\[
J(\beta+v)=J(\beta).
\]
Hence minimizers are not unique: infinitely many $\beta$ give the same fit.

Geometrically, moving along null directions does not change projection $X\beta$ in data space.

Ridge solves
\[
\hat\beta_\lambda=\argmin_\beta\|y-X\beta\|_2^2+\lambda\|\beta\|_2^2,
\]
with closed form
\[
\hat\beta_\lambda=(X^\top X+\lambda I)^{-1}X^\top y.
\]
For $\lambda>0$, $X^\top X+\lambda I\succ0$, so solution is unique.
As $\lambda\to0^+$, $\hat\beta_\lambda$ approaches the minimum-norm OLS solution $X^+y$ (Moore--Penrose pseudoinverse solution).

\subsection*{Q3. High $R^2$ but Bad Test Error}
Example construction:
\begin{itemize}
\item Training data: $x_i\in[-1,1]$, $y_i=\sin(\pi x_i)+\epsilon_i$ with small noise.
\item Fit a very high-degree polynomial (e.g., degree $n-1$ interpolation).
\end{itemize}
Training $R^2$ can be extremely high (often near $1$), but test error is poor because the learned curve has high oscillation between training points (high variance).

Bias--variance view at a test point $x$:
\[
\E\big[(Y-\hat f(x))^2\big]
=\underbrace{\sigma^2}_{\text{noise}}+\underbrace{\big(\E[\hat f(x)]-f(x)\big)^2}_{\text{bias}^2}+\underbrace{\operatorname{Var}(\hat f(x))}_{\text{variance}}.
\]
Over-parameterized models can have low bias and high variance, so generalization degrades despite high training $R^2$.

\subsection*{Q4. Hat Matrix and Influence}
For full-column-rank $X\in\R^{n\times p}$,
\[
H=X(X^\top X)^{-1}X^\top.
\]
We need $\sum_i h_{ii}=p$.

Use trace:
\[
\sum_{i=1}^n h_{ii}=\operatorname{tr}(H)
=\operatorname{tr}\big(X(X^\top X)^{-1}X^\top\big)
=\operatorname{tr}\big((X^\top X)^{-1}X^\top X\big)
=\operatorname{tr}(I_p)=p.
\]

Interpretation: total leverage equals number of fitted parameters. So $\operatorname{tr}(H)$ behaves like model complexity (effective degrees of freedom).

Why high leverage points distort fit:
\begin{itemize}
\item $\hat y=Hy$, so $\hat y_i=\sum_j h_{ij}y_j$.
\item Large $h_{ii}$ means point $i$ strongly controls its own fitted value.
\item If such a point is extreme in feature space (and/or has large residual), it can rotate/shift the regression hyperplane disproportionately.
\end{itemize}

\section{Logistic Regression}

\subsection*{Q5. Convexity of Logistic Loss}
Binary labels $y_i\in\{0,1\}$, scores $z_i=\xx_i^\top\beta$, probabilities $p_i=\sigma(z_i)=1/(1+e^{-z_i})$.
Negative log-likelihood:
\[
\ell(\beta)=\sum_{i=1}^n\left[\log(1+e^{z_i})-y_i z_i\right].
\]
Gradient:
\[
\nabla \ell(\beta)=X^\top(p-y).
\]
Hessian:
\[
\nabla^2\ell(\beta)=X^\top W X,
\quad W=\operatorname{diag}(p_i(1-p_i)).
\]
Each $p_i(1-p_i)\ge0$, so for any $v$,
\[
v^\top\nabla^2\ell(\beta)v=(Xv)^\top W(Xv)=\sum_i w_i(Xv)_i^2\ge0.
\]
Hence Hessian is PSD and logistic loss is convex.

Convexity implies every local minimum is global. If Hessian is PD on feasible directions (e.g., full rank plus no complete separation), minimizer is unique.

\subsection*{Q6. Decision Boundary and Feature Maps}
Logistic model:
\[
\Pp(y=1\mid x)=\sigma(\ww^\top x+b).
\]
Decision rule at threshold $0.5$:
\[
\sigma(\ww^\top x+b)\ge0.5\iff \ww^\top x+b\ge0.
\]
So boundary is linear in original features.

If we use transformed features $\phi(x)$,
\[
\Pp(y=1\mid x)=\sigma(\ww^\top\phi(x)+b).
\]
Boundary is linear in $\phi$-space but can be nonlinear in original $x$-space. So logistic regression remains a linear classifier in the chosen feature representation.

\subsection*{Q7. Extreme Class Imbalance (99\% class 0)}
\textbf{Why accuracy misleads:}
Predicting always class $0$ gives 99\% accuracy but 0 recall for class $1$.

\textbf{Behavior of logistic regression:}
MLE emphasizes likelihood across all points; with severe imbalance, intercept tends strongly negative and predicted probabilities often stay below 0.5.

\textbf{Decision threshold effect:}
Default threshold 0.5 usually yields many false negatives. In deployment one should tune threshold by costs/ROC/PR curve.

\textbf{Optimal threshold under unequal costs:}
Let $\eta(x)=\Pp(Y=1\mid x)$.
\begin{itemize}
\item Cost of predicting 1: $C_{FP}\Pp(Y=0\mid x)=C_{FP}(1-\eta)$.
\item Cost of predicting 0: $C_{FN}\Pp(Y=1\mid x)=C_{FN}\eta$.
\end{itemize}
Predict class 1 when
\[
C_{FP}(1-\eta)\le C_{FN}\eta
\iff
\eta\ge\frac{C_{FP}}{C_{FP}+C_{FN}}.
\]
Equivalent likelihood-ratio form:
\[
\frac{p(x\mid1)}{p(x\mid0)}\ge\frac{C_{FP}}{C_{FN}}\frac{\pi_0}{\pi_1}.
\]

\section{Perceptron}

\subsection*{Q8. Non-Separable Data}
Perceptron update on mistakes:
\[
\ww_{t+1}=\ww_t+y_i\xx_i.
\]
If algorithm stopped after finite steps, then no future mistakes occur, so final hyperplane correctly classifies all training points. That implies separability, contradiction. Therefore for non-separable data, perceptron keeps making mistakes forever.

Asymptotic behavior: weights generally do not converge to a fixed point; they keep moving among conflicting constraints, often with growing norm and oscillating direction.

Soft-margin SVM comparison:
\begin{itemize}
\item Perceptron optimizes no regularized convex objective in non-separable case.
\item Soft-margin SVM solves convex optimization with slack, always well-posed and convergent.
\item SVM returns a controlled margin/error tradeoff; perceptron does not.
\end{itemize}

\subsection*{Q9. Mistake Bound and Scaling}
Classical bound for separable data:
\[
T\le\left(\frac{R}{\gamma}\right)^2,
\]
where $R=\max_i\|\xx_i\|$ and $\gamma$ is geometric margin. This is equivalent to $T\le(RB)^2$ with $B=1/\gamma$.

Effects of scaling:
\begin{itemize}
\item Uniform scaling $\xx\mapsto a\xx$ scales both $R$ and $\gamma$ by $a$, so ratio $R/\gamma$ unchanged.
\item Non-uniform scaling can increase $R$ and shrink effective margin, worsening bound.
\item Feature normalization tends to control $R$ and improve practical convergence speed.
\end{itemize}

Can bad scaling cause divergence?
\begin{itemize}
\item If data is linearly separable, perceptron still converges theoretically (finite mistakes), though may be extremely slow numerically.
\item For non-separable data, it does not converge regardless of scaling.
\end{itemize}

\section{Support Vector Machines}

\subsection*{Q10. Hard vs Soft Margin Limits}
Soft-margin primal:
\[
\min_{\ww,b,\xi\ge0}\frac12\|\ww\|^2+C\sum_{i=1}^n\xi_i,
\quad y_i(\ww^\top x_i+b)\ge1-\xi_i.
\]

If data is separable, there exist feasible points with $\xi_i=0$. As $C\to\infty$, any nonzero slack becomes very expensive, so optimizer approaches zero-slack solution, i.e. hard-margin SVM.

As $C\to0$, slack is nearly free, objective mainly minimizes $\|\ww\|^2$, so $\ww\to0$ (trivial classifier tendency), with many margin violations tolerated.

\subsection*{Q11. KKT Interpretation for Soft-Margin SVM}
Dual constraints: $0\le\alpha_i\le C$.
KKT:
\[
\alpha_i\big(y_if_i-1+\xi_i\big)=0,
\quad (C-\alpha_i)\xi_i=0,
\quad f_i=\ww^\top x_i+b.
\]

Cases:
\begin{enumerate}
\item $\alpha_i=0$: then $\xi_i=0$ and $y_if_i\ge1$ (typically outside margin). Not a support vector.
\item $0<\alpha_i<C$: then $\xi_i=0$ and $y_if_i=1$. Exactly on margin (support vector).
\item $\alpha_i=C$: on/inside margin. Usually $y_if_i\le1$; if $y_if_i<1$ then inside margin, and if $y_if_i<0$ point is misclassified.
\end{enumerate}

\subsection*{Q12. Margin Geometry}
Canonical separating hyperplanes:
\[
\ww^\top x+b=1\quad\text{and}\quad \ww^\top x+b=-1.
\]
Distance between parallel hyperplanes $\ww^\top x+b=c_1$ and $c_2$ is
\[
\frac{|c_1-c_2|}{\|\ww\|}.
\]
Hence margin width is
\[
\frac{|1-(-1)|}{\|\ww\|}=\frac{2}{\|\ww\|}.
\]
So maximizing margin is equivalent to minimizing $\|\ww\|$, or equivalently minimizing $\frac12\|\ww\|^2$.

\subsection*{Q13. Why the Dual Uses Inner Products; PSD Requirement}
Soft/hard SVM dual objective has form
\[
\max_\alpha\sum_i\alpha_i-\frac12\sum_{i,j}\alpha_i\alpha_j y_i y_j\langle x_i,x_j\rangle,
\]
so data appears only through pairwise inner products. Replacing
$\langle x_i,x_j\rangle$ by $K(x_i,x_j)$ gives kernel SVM.

If $K$ is not PSD, Gram matrix can be indefinite. Then the dual is not guaranteed concave (optimization may become non-convex with local optima), and no valid RKHS/Mercer interpretation is guaranteed.

\section{KNN and Bayes}

\subsection*{Q14. Curse of Dimensionality Distance Concentration}
Let $X,Y\sim\text{Unif}([0,1]^d)$ independent. Define
\[
D^2=\|X-Y\|_2^2=\sum_{j=1}^d (X_j-Y_j)^2.
\]
Each term is bounded and i.i.d. with mean $1/6$. By law of large numbers,
\[
\frac{D^2}{d}\to\frac16\quad\text{(in probability)}.
\]
So pairwise distances concentrate around $\sqrt{d/6}$. Spread is order $\sqrt d$, mean is order $d$ in squared distance, so relative spread shrinks as $1/\sqrt d$.

Therefore nearest and farthest neighbor distances become similar in relative terms, i.e.
\[
\frac{D_{\max}-D_{\min}}{D_{\min}}\to0,
\]
which breaks ``nearest means similar'' intuition and hurts KNN.

\subsection*{Q15. Bayes Error vs 1-NN; Bayes Consistency}
At a point $x$, let $\eta(x)=\Pp(Y=1\mid X=x)$.

Bayes conditional error:
\[
R^*(x)=\min\{\eta(x),1-\eta(x)\}.
\]
As sample size grows, the nearest neighbor label behaves like an independent Bernoulli draw with parameter $\eta(x)$. Then 1-NN conditional error tends to
\[
R_{1NN}(x)=2\eta(x)(1-\eta(x)).
\]
Since $2\eta(1-\eta)\le2\min\{\eta,1-\eta\}$,
\[
R_{1NN}(x)\le 2R^*(x).
\]
This gives the standard bound intuition.

Why $k\to\infty$ with $k/n\to0$ is Bayes-consistent:
\begin{itemize}
\item $k\to\infty$: enough neighbors to average labels, reducing variance.
\item $k/n\to0$: neighbor radius shrinks to $0$, so local class probability estimate is asymptotically unbiased.
\end{itemize}
Hence risk converges to Bayes risk.

\subsection*{Q16. Scaling Trap Construction}
Let training set (for 1-NN):
\[
(0,1000)\to0,\ (0,900)\to0,\ (1,0)\to1,\ (1,100)\to1.
\]
Query: $q=(1,950)$ (informative feature is first coordinate, so true class should be 1).

Without normalization:
\[
d(q,(0,900))=\sqrt{1+50^2}\approx50.01,
\]
\[
d(q,(1,100))=850,
\]
so nearest neighbor is class 0 (wrong).

After min-max normalization per feature:
\[
x_1\in[0,1],\quad x_2\in[0,1000],
\]
query becomes $(1,0.95)$, and
\[
d((1,0.95),(1,0.1))=0.85,
\]
\[
d((1,0.95),(0,0.9))\approx1.001,
\]
so nearest is class 1 (correct). This shows distance domination by large-scale irrelevant features.

\section{Decision Trees and Random Forest}

\subsection*{Q17. Perfect Feature and Information Gain}
Information gain for feature $X$:
\[
IG(Y,X)=H(Y)-H(Y\mid X).
\]
If $X$ perfectly determines $Y$, then each split child is pure, so
\[
H(Y\mid X)=0.
\]
Therefore
\[
IG(Y,X)=H(Y).
\]
So gain equals initial entropy.

\subsection*{Q18. Why Unrestricted Trees Have Infinite VC Dimension}
Take any $m$ distinct real points $x_1<\cdots<x_m$ on $\R$.
A deep enough axis-aligned tree can place split thresholds between consecutive points to isolate each point in its own leaf. Then assign each leaf label 0 or 1 arbitrarily. So every labeling of these $m$ points is realizable.
Since $m$ is arbitrary, VC dimension is infinite.

\subsection*{Q19. Why Bagging Reduces Variance but Not Bias}
Let $\hat f_b(x)$ be predictor from bootstrap sample $b$, and ensemble
\[
\bar f_M(x)=\frac1M\sum_{b=1}^M\hat f_b(x).
\]
For squared loss:
\[
\text{MSE}=\underbrace{\big(\E[\bar f_M]-f\big)^2}_{\text{bias}^2}+\underbrace{\operatorname{Var}(\bar f_M)}_{\text{variance}}+\sigma^2.
\]
Averaging mostly affects variance. If base learners have variance $\sigma_f^2$ and pairwise correlation $\rho$,
\[
\operatorname{Var}(\bar f_M)=\rho\sigma_f^2+\frac{1-\rho}{M}\sigma_f^2.
\]
As $M$ grows, second term vanishes, first remains. Bias comes from model class and is not fundamentally removed by averaging.

\subsection*{Q20. Why Small Feature Subset Size in RF Reduces Correlation}
If all trees can always use strongest predictors, many trees split similarly, increasing correlation $\rho$.
Random feature subsampling forces diverse split choices, decreasing $\rho$ in
\[
\operatorname{Var}(\bar f_M)=\rho\sigma_f^2+\frac{1-\rho}{M}\sigma_f^2.
\]
Lower $\rho$ lowers ensemble variance substantially. Tradeoff: very small subset size may weaken each tree (slightly higher bias).

\section{Model Evaluation}

\subsection*{Q21. Accuracy Paradox Example}
Confusion matrix (1000 samples, 5\% positives):
\[
\begin{array}{c|cc}
&\hat y=1&\hat y=0\\\hline
y=1&0&50\\
y=0&0&950
\end{array}
\]
Accuracy $=(950+0)/1000=95\%$.
Model is useless for positive class (recall $=0$).

\subsection*{Q22. Higher Accuracy but Lower F1}
Assume 1000 samples, 100 positives.

Model A:
\[
TP=10,\ FP=0,\ FN=90,\ TN=900.
\]
Accuracy $=91\%$,
\[
F1_A=\frac{2\cdot1\cdot0.1}{1+0.1}=0.1818.
\]

Model B:
\[
TP=70,\ FP=120,\ FN=30,\ TN=780.
\]
Accuracy $=85\%$,
\[
\text{Precision}=\frac{70}{190}=0.3684,\quad
\text{Recall}=0.7,
\quad F1_B\approx0.4828.
\]
So B has lower accuracy but much better F1 and is preferable when catching positives matters.

\subsection*{Q23. LOOCV and High Leverage}
Given
\[
\hat y_i^{(-i)}=\frac{\hat y_i-h_{ii}y_i}{1-h_{ii}}.
\]
Then LOOCV residual is
\[
e_i^{(-i)}=y_i-\hat y_i^{(-i)}
=\frac{y_i-\hat y_i}{1-h_{ii}}=\frac{e_i}{1-h_{ii}}.
\]
So LOOCV squared error term scales as
\[
\big(e_i^{(-i)}\big)^2=\frac{e_i^2}{(1-h_{ii})^2}.
\]
If $h_{ii}\approx1$, denominator is tiny and contribution explodes. This is why highly leveraged points can dominate LOOCV error.

\section{Multiclass Classification}

\subsection*{Q24. One-vs-All Failure Case}
Construct 3-class data in $\R^2$:
\begin{itemize}
\item Class A near $(-2,0)$, many points.
\item Class B near $(2,0)$, many points.
\item Class C near $(0,0)$, very few points.
\end{itemize}
For OvA classifier of class C (C vs A+B), negatives dominate and are on both sides of C; with linear model and imbalance, optimal score for C can be negative almost everywhere. Then argmax over OvA scores never predicts C, so the class disappears.

Why joint multiclass SVM helps: it learns all class scores together with constraints
\[
\ww_{y_i}^\top x_i-\ww_j^\top x_i\ge1-\xi_i,\quad \forall j\neq y_i,
\]
forcing C examples to beat all other classes by margin. This coupling prevents trivial elimination of a class.

\subsection*{Q25. All-Pairs (OvO) Complexity}
For $k$ classes:
\begin{itemize}
\item OvA: train $k$ classifiers.
\item OvO: train $k(k-1)/2$ classifiers.
\end{itemize}
So training/model storage is $O(k^2)$ for OvO.

When infeasible: large $k$ (e.g., $k=1000$ gives $499{,}500$ classifiers; $k=10{,}000$ gives about $50$ million).

When OvO may be better than OvA:
\begin{itemize}
\item each binary problem is smaller and often better balanced,
\item pairwise boundaries may be simpler,
\item can improve accuracy despite more models.
\end{itemize}
Prediction cost is also heavier for OvO (often $O(k^2)$ voting, unless structured decoding like DAG is used).

\section{Very Advanced}

\subsection*{Q26. Bias--Variance for Regression; Why Classification is Harder}
Let
\[
Y=f(x)+\varepsilon,\quad \E[\varepsilon\mid x]=0,\quad \operatorname{Var}(\varepsilon\mid x)=\sigma^2.
\]
For estimator $\hat f_D(x)$ (random over dataset $D$):
\[
\E_{D,\varepsilon}\big[(Y-\hat f_D(x))^2\mid x\big]
=\sigma^2+\big(\E_D[\hat f_D(x)]-f(x)\big)^2+\operatorname{Var}_D(\hat f_D(x)).
\]
This is exact decomposition for squared loss.

For classification with 0--1 loss, risk uses indicator
$\1\{\hat y\neq y\}$, which is discontinuous/non-quadratic. There is no equally clean additive bias+variance decomposition in general. One typically uses surrogate losses, margin analysis, or approximate decompositions.

\subsection*{Q27. VC Dimension Comparison}
\begin{itemize}
\item Linear classifiers in $\R^d$: $\mathrm{VC}=d+1$.
\item Decision trees of depth $k$ (axis-aligned, in $\R^d$): capacity grows exponentially with depth; typical bounds are on the order of $O(2^k\log d)$ (and at least $\Omega(2^k)$). Unrestricted depth gives infinite VC dimension.
\item SVM with RBF kernel: without norm regularization, effective VC dimension is infinite (can shatter arbitrarily many points). Regularization controls effective capacity in practice.
\end{itemize}

\subsection*{Q28. Perceptron vs Logistic vs Hard SVM}
\begin{center}
\begin{tabular}{>{\raggedright\arraybackslash}p{0.18\linewidth} >{\raggedright\arraybackslash}p{0.25\linewidth} >{\raggedright\arraybackslash}p{0.22\linewidth} >{\raggedright\arraybackslash}p{0.27\linewidth}}
\toprule
Method & Loss / Objective & Margin behavior & Convergence / Robustness \\
\midrule
Perceptron & Minimize mistakes implicitly via updates on misclassified points (non-convex objective surrogate) & No explicit max-margin objective; in separable case finds some separator & Converges only for separable data; unstable on non-separable; sensitive to scaling/outliers \\
\addlinespace
Logistic regression & Convex log-loss $\sum\log(1+e^{-y f(x)})$ (often with L2/L1 regularization) & Encourages large positive margins smoothly but not max-margin exactly & Global optimum (convex); probabilistic outputs; more robust optimization than perceptron \\
\addlinespace
Hard-margin SVM & $\min \frac12\|w\|^2$ s.t. $y_if_i\ge1$ & Explicitly maximizes geometric margin & Convex QP with unique $w$ for separable data; sensitive to label noise/outliers (no slack) \\
\bottomrule
\end{tabular}
\end{center}

\paragraph{Bottom-line comparison.}
Perceptron is simplest but least robust outside separability. Logistic is probabilistic and convex. Hard SVM gives strongest geometric margin but requires clean separability; soft-margin SVM is the practical compromise.

\end{document}
